\section{Introduction}
\label{sec:introduction}

Rust prevents broad classes of memory errors through its ownership and borrowing disciplines and has gained increasing adoption in systems software~\cite{jung2017rustbelt,li2024empirical}.
However, memory safety alone does not guarantee functional correctness or security~\cite{xu2021memory}. A memory-safe Rust program may still overflow, violate a protocol invariant, or produce an incorrect result~\cite{Edera}.
Formal verification provides a rigorous way to establish such functional-correctness properties by proving that a program satisfies a precise specification.

% Rust formal verification, Verus, Prusti Cresuot ... ATP,  for ITP, Charon-Aeneas LLBC benifits xxx supports xxx, currently not Isabelle/HOL.

% we select Rust2isabelle/HOL. Isabelle benifits xxx, 

% Cresuot-why3-Isabelle/HOL (why3 common experssion), Autocorrede (manually model in Isabelle/HOL). We select charon-Aeneas path xxx for the reason xxx.

Isabelle/HOL~\cite{nipkow2002isabelle} is a mature interactive theorem prover based on classical higher-order logic.
It supports recursive definitions, algebraic datatypes, reusable theorem libraries, and substantial proof automation~\cite{bohme2010sledgehammer,nelson2019scaling}.
In addition, it has supported several large-scale systems verification efforts~\cite{alkassar2008formal,becker2026nitro,zhao2019rely}, most notably the functional-correctness verification of the seL4 microkernel~\cite{sel4-sosp}.
These features make Isabelle/HOL a natural target for reasoning about the functional correctness of Rust programs.


However, manually constructing and maintaining a faithful Isabelle/HOL model of Rust programs requires considerable effort. 
% Such a model must account for Rust's rich type system. 
% A direct translation must handle ownership and borrowing without reintroducing unnecessary memory reasoning.
In particular, a direct translation need preserve Rust's ownership and borrowing discipline without reintroducing an explicit memory model.
The Charon--Aeneas~\cite{ho2025charon,ho2022aeneas} toolchain provides a reusable foundation for automating this process.
% We therefore build on the relatively mature Charon--Aeneas toolchain rather than developing the entire translation infrastructure 
% About the toolchain, 
Charon translates Rust compiler Mid-level Intermediate Representation (MIR) into the Low-Level Borrow Calculus (LLBC), a typed and structured intermediate representation designed for downstream program analysis.
Aeneas then consumes LLBC and uses symbolic execution to eliminate ownership and borrowing constructs, producing a Pure functional AST suitable for verification.
Aeneas currently provides extraction backends for F$^\star$~\cite{swamy2016dependent}, Rocq~\cite{bertot2013interactive}, HOL4~\cite{slind2008brief}, and Lean4~\cite{moura2021lean}, but it lacks support for Isabelle/HOL.

To fill this gap, we develop \texttt{Rust2Isabelle} as a new extraction backend for Aeneas.
It combines automatic generation of Isabelle/HOL theories with a companion library defining the Rust primitives used in the output.

We provide the following contributions.
\begin{itemize}
    \item \textbf{Rust2Isabelle Translation Rules}: We define translation rules from Rust to Isabelle/HOL. 
    These rules cover several Rust features, including types, expressions, mutable borrows, traits, const generics, and control flow.
    \item \textbf{Aeneas-Isabelle Prototype}: We implement a new Isabelle/HOL backend for Aeneas. 
    This backend generates complete Isabelle/HOL theory files from Rust code. 
    Our implementation is written in OCaml and integrates easily into the existing Aeneas ecosystem.
    \item \textbf{Evaluation}: We evaluate our approach and implementation using the official Aeneas test suite. 
    Additionally, we apply our tool to a key component of the real-world Solana~\cite{SolanaRBPF2} blockchain virtual machine as a case study.
\end{itemize}

% 1. Translation rules from Rust 2 Isabelle
% 2. compiler tool implementation + intergration
% 3. evaluation

{
\def\sectionautorefname{Section}
The rest of the paper is organized as follows.
\autoref{sec:background} reviews the relevant features of Rust, Aeneas's Pure AST, and Isabelle/HOL. 
\autoref{sec:design} presents the \texttt{Rust2Isabelle} workflow and the translation rules by language construct.
\autoref{sec:evaluation} combines construct-level coverage with an SBPF (Solana eBPF) virtual machine case study, Isabelle/HOL acceptance, and selected semantic properties.
Related work and the remaining limitations follow in~\autoref{sec:related} and~\autoref{sec:conclusion}.
}